% [原文 p.47]
\chapter{极小模型纲领引论}

本章初次接触一般的极小模型纲领，并处理若干预备性题材。

§2.1 不拘泥于形式，说明极小模型纲领的目的与方法。一个基本洞察是：三维以上，极小模型纲领不会停留在非奇异簇的类别之内。因此，若要精确解说极小模型纲领，就必须首先处理必然出现的奇点。

§2.2 稍偏离“极小模型纲领入门”这一主线，考察极小模型纲领的若干推广。这些推广在应用时极为有益，却不要求任何困难的概念。

最有益的大概是所谓的对数（log）范畴的研究。其中处理簇 $X$ 与素除子的形式线性组合 $D$ 组成的对 $(X,D)$。对极小模型纲领而言，这似乎会成为自然的设定。

这类对 $(X,D)$ 的奇点的若干类在 §2.3 中考察。它们稍带技术性，却是后续发展不可缺少的。对数范畴更详细的处理见 [Kol97]。

§2.4、§2.5 的主题是证明后续各章要用的消没定理。我们只打算证明日后所需的最小限度，故把讨论限制在非奇异射影簇的情形。

% [原文 p.48]
此情形下的证明不仅简单，还揭示了消没定理与簇的拓扑之间的关系。消没定理的证明已有若干种。其他处理方式见 [KMM87, EV92, Kol95, Kol97] 等。

\section{森纲领入门}

\takeaway{本节从“为什么必须允许奇点”出发：先用例 2.1 说明三维以上非奇异簇不够用，再给出锥定理、极射线收缩与翻转的框架，定义标准奇点、终端奇点与极小模型，最后以四步流程勾画森纲领。}

如定理 1.32 之后所述，在 E3、E4、E5 的情形下，收缩所得簇 $Y$ 上出现的奇点使定理 1.32 无法适用。因此，必须确认当 $X$ 具有某种奇点时，定理 1.32 的类似结论能否成立。实际上，E3 与 E4 情形下的奇点并不构成大问题。更一般地，当 $X$ 只有孤立超曲面奇点时，证明的大部分只需稍作修补（参见 [Cut88a]）。

然而，E5 情形下的奇点不是超曲面奇点，直接对具有这类奇点的簇作几何研究要困难得多。甚至在 [Mor82] 之前，人们就知道这类奇点会带来麻烦。例如，[Uen75] 作了如下考察。

\begin{example}{2.1}
设 $A$ 为三维 Abel 簇，$X$ 为按二阶自同构 $x\mapsto -x$ 取 $A$ 的商所得的空间 $A/\pm$。
2-挠点被该自同构固定，故它们给出 $X$ 上 64 个与 $\C^3/\pm$ 解析同构的奇点。[Uen75, 16.17] 证明了具有 nef 标准类的非奇异射影簇不与 $X$ 双有理等价。这意味着只要停留在非奇异簇的范围内，定理 1.29 在三维就不存在类似结论。
\end{example}
% [Tutor Note: 例 2.1 说明为何不能只在非奇异簇中工作。该商奇点是 $\C^3/\pm$ 型的终端奇点，任何双有理模型都无法把它消去；故三维极小模型纲领必须改换容许奇点的类。]

\par\medskip\noindent\textbf{2.2}\quad 需要一种引入新概念的研究方法。许多问题都是如此，这里也存在两种相反的立场。
\begin{itemize}[leftmargin=2.2em,itemsep=2pt]
\item 最小主义：从非奇异簇出发，把应用定理 1.32 或其推广时出现的奇点类中最小者确立为概念。
\item 最大主义：把定理 1.29 的类似结论可能成立的奇点类中最大者确立为概念。
\end{itemize}

% [原文 p.49]
本节从最大主义的方针开始。因为要展开一般理论，这一方针最为合适。但若不对奇点加上更强的条件，就得不出像样的结果。因此，到了证明的阶段，就不得不采取接近最小主义的态度。\hfill$\square$

\par\medskip\noindent\textbf{2.3}\quad 在概要 1.31(1) 中，首先考察 $K_X$ 是否为 nef。这一提问要有意义，就必须定义 $K_X$ 与曲线之间的交点数。交点数在除子类 $K_X$ 为Cartier 时、更一般地在 $K_X$ 为 $\Q$-Cartier 时才有定义。看来三维以上若没有这一假设，本质上得不到任何东西。

接下来必须对 $X$ 建立锥定理。对非奇异簇，这已由定理 1.24 完成。非奇异性之所以必要，是因为用到了定理 1.10 步骤 2 中的公式；当 $X$ 有奇点时，该公式往往不成立。然而存在使定理 1.10 本身成立的更大的奇点类，这引出了许多带奇点簇的锥定理。事实上，锥定理对任意三维正规簇成立 [Kol92, p.295]，但在更高维不成立 [Kol92, 5.5.2.2.1]。足以在四维以上执行极小模型纲领的锥定理，目前用几何方法尚未证明。用不同方法证明锥定理将在第 3 章给出。\hfill$\square$

\par\medskip\noindent\textbf{2.4}\quad 在概要 1.31 步骤 2 中，需要收缩极射线 $R$。这似乎是更精细的结果，看来只对相当特殊的奇点成立；但第 3 章给出的证明次序恰好相反：先证明极射线的收缩，然后才证明锥定理。\hfill$\square$

极射线的收缩分为三种类型，首先叙述之。后面的讨论将表明，这三种情形具有完全不同的特征。

\begin{propositionp}{2.5}
设 $X$ 为正规射影簇，并假设 $X$ 是 $\Q$-分解的（$\Q$-factorial）。若 $f\colon X\to Y$ 是极射线 $R\subset\NE(X)$ 的收缩，则下列之一成立。
\begin{enumerate}[label=（\arabic*）,leftmargin=2.6em,itemsep=1pt]
\item （纤维收缩）$\dim X>\dim Y$。
\item （除子收缩）$f$ 为双有理的，且 $\Exc(f)$ 是不可约除子。
% [原文 p.50]
\item （小收缩）$f$ 为双有理的，且 $\Exc(f)$ 的余维数 $\ge 2$。
\end{enumerate}
\end{propositionp}

% [Tutor Note: 三种收缩与概要 1.31 的子步骤对应：纤维收缩对应步骤 4，除子收缩对应步骤 2，小收缩则是步骤 3 才出现的新现象，且只可能在三维以上发生。]

\begin{proof}
需要证明的只是如下断言：若 $f$ 为双有理的，且例外集含有不可约除子 $E$，则 $E=\Exc(f)$。由 $E$ 为 $\Q$-Cartier 这一假设可知 $(E\cdot R)<0$。这将在引理 3.39 中证明。若 $C\subset X$ 是使 $f(C)$ 为一点的曲线，则 $[C]\in R$，故 $(E\cdot C)<0$。这特别意味着 $C\subset E$，从而 $E=\Exc(f)$。
\end{proof}

\par\medskip\noindent\textbf{2.6}\quad 命题 2.5 的三种情形，从我们的纲领的立场看，行为大为不同。

\par\medskip\noindent\textbf{情形 1}\quad 纤维收缩对应于概要 1.31 步骤 4。若 $(R\cdot K_X)<0$，则 $f$ 的一般纤维 $F$ 是使反标准除子 $-K_F=-K_X|_F$ 为丰富的代数簇。因此，至少在原则上，把理解 $X$ 的问题归结为对底空间 $Y$ 与 $f$ 的纤维等低维簇的研究。此外，纤维是 $\C\PP^1$ 与 Del Pezzo 曲面之类的类似物，是极为特殊的簇。

\par\medskip\noindent\textbf{情形 2}\quad 除子收缩应当对应于概要 1.31 步骤 2。要把这一点说到底，必须证明至少 $K_Y$ 为 $\Q$-Cartier ，或者更强地，$Y$ 是 $\Q$-分解的。若后者成立，则 $\rho(Y)\le\rho(X)-1$，故可以认为 $Y$ 比 $X$“更简单”。

设 $D\subset Y$ 为任意Weil 除子，$D'\subset X$ 为其双有理变换。适当选取有理数 $a(D)$，使 $(R\cdot(D'+a(D)E))=0$ 成立。最后选取自然数 $m$，使 $D'':=m(D'+a(D)E)$ 为Cartier 除子。如此定义的 $D''$ 很可能成为 $Y$ 的Cartier 除子 $\bar D$ 的拉回。假如确实如此，则 $mD\sim\bar D$，故 $D$ 为 $\Q$-Cartier 。

\par\medskip\noindent\textbf{情形 3}\quad 小收缩。这是新的状况。这种状况在曲面情形下（由于维数的关系）绝不会发生，在非奇异三维簇的情形下也未曾发生。在小收缩的情形下，$Y$ 具有“非常坏”的奇点：标准除子 $K_Y$ 无论乘多少倍都不是Cartier 的。

% [原文 p.51]
实际上，试设 $mK_Y$ 与 $mK_X$ 都是Cartier 的。则 $mK_X$ 与 $f^*(mK_Y)$ 是 $X$ 的Cartier 除子，且在余维数 2 的子集之外线性等价。因此它们在 $X$ 上整体线性等价。然而由 $(R\cdot K_X)<0$ 与 $(R\cdot f^*(mK_Y))=0$，这是不可能的。

这样一来，我们就走出了可控制的簇类。此时要继续纲领，必须引入称为翻转（flip）的新操作。这一操作是拓扑手术（surgery）的代数几何类似物。

其做法是：不收缩余维数 $\ge 2$ 的子簇 $E=\Exc(f)\subset X$，而先将其除去，然后再添加同样余维数 $\ge 2$ 的子簇 $E^+$，把 $X\setminus E$ 紧化。（眼下，这一操作是否存在、即使它改善状况是否还有准确的定义，都还什么都不清楚，暂且搁置。）\hfill$\square$

\begin{example}{2.7}
（翻转的例子）用例子考察这一状况。在此例中，翻转把曲线 $C\cong\C\PP^1$ 从奇异簇 $X_n$ 中除去，代之以 $D\cong\C\PP^1$，得到“改良的”簇 $X^+_n$。在此情形下，它是非奇异的。

先从辅助的构造开始。考察簇 $Y=(xy-uv=0)\subset\C^4$。它在原点有孤立奇点。将它爆破，得到 $\bar X=B_0Y$。例外集 $Q\subset\bar X$ 同构于射影二次曲面 $(xy-uv=0)\subset\PP^3$。它含有两族直线：一族由 $x=cv,\ y=c^{-1}u$ 给出，另一族由 $x=cu,\ y=c^{-1}v$ 给出（$c$ 为参数）。把第一族中的各直线收缩，得到非奇异簇 $X$。由第二族直线同样得到非奇异簇 $X^+$。$X$（或 $X^+$）也可以由爆破 $Y$ 中理想 $(x,v)$（或 $(x,u)$）得到。设 $C\subset X$（或 $C^+\subset X^+$）为射 $X\to Y$（或 $X^+\to Y$）的例外曲线。此时得到如下簇与射的图式。

\begin{center}
\begin{tikzcd}[column sep=3.2em, row sep=2.8em]
{} & \bar X \arrow[rr] \arrow[dl] \arrow[dr] & {} & \bar X_n \arrow[dll] \arrow[d] \\
X \arrow[r] \arrow[dr] & X_n \arrow[dr] & X^{+} \arrow[r] \arrow[dl] & X^{+}_n \arrow[dl] \\
{} & Y \arrow[r] & Y_n & {}
\end{tikzcd}
\end{center}

% [原文 p.52]
考虑 $n$ 阶循环群的作用 $\mu_n\colon(x,y,u,v)\mapsto(\zeta x,y,\zeta u,v)$。由此，上述所有簇上都定有作用。相应的商簇用下标 $n$ 表示。

在 $Y$ 上，作用的固定点集（即被某个 $g\in\mu_n\setminus\{1\}$ 固定的点的全体）是二维平面 $(x=u=0)$。在射影二次曲面 $Q$ 上，固定点集是如下两条直线：
对应于固定平面的 $(x=u=0)$，与对应于 $\zeta$-特征空间的 $(y=v=0)$。
因此 $X$ 上的固定点集由如下两个连通分支组成：
平面 $(x=u=0)$ 的双有理变换，与直线 $(y=v=0)$ 的像。
后者是孤立固定点。容易看出，$(x,v'=vx^{-1},u)$ 给出该孤立固定点处的局部坐标，群作用由 $(x,v',u)\mapsto(\zeta x,\zeta^{-1}v',\zeta u)$ 给出。

在 $X^+$ 上，固定点集是不可约除子，且含有例外曲线 $C^+$。因此 $X^+_n$ 是非奇异的。

\begin{figure}[htbp]
\centering
\resizebox{0.72\textwidth}{!}{\begin{tikzpicture}[line width=0.5pt]
% 左图：X
\fill[gray!30] (0.00,3.20) -- (0.00,0.80) -- (1.15,2.00) -- cycle;
\draw (0.00,3.20) -- (1.15,2.00) -- (0.00,0.80) -- cycle;
\draw (1.15,2.00) -- (2.05,2.00);
\fill (2.05,2.00) circle (0.075);
\draw (2.05,2.00) -- (3.30,3.30);
\draw (2.05,2.00) -- (3.10,1.05);
\node at (0.60,3.70) {$X$};
% 右图：X^{+}
\fill[gray!30] (5.30,3.30) -- (6.15,2.00) -- (5.30,0.70) -- cycle;
\draw (5.30,3.30) -- (6.15,2.00) -- (5.30,0.70) -- cycle;
\draw (6.15,2.00) -- (7.35,3.30);
\draw (6.15,2.00) -- (7.35,0.70);
\node at (5.90,3.95) {$X^{+}$};
% 翻模
\node at (4.45,2.98) {翻模};
\draw[dashed,->] (3.92,2.55) -- (4.98,2.55);
% 群作用的不动点集
\node[align=center] (fix) at (3.65,0.55) {群作用的\\不动点集};
\draw (fix.north west) -- (0.52,0.72);
\draw (fix.north east) -- (5.34,1.60);
\end{tikzpicture}
}
\figtitle{2.1}{翻模按群取商诱导翻转。固定点集中余维数 $\ge 2$ 的部分产生奇点，余维数 1 的部分产生 $K$ 的符号（原书 p.52）。}
\end{figure}
% [Tutor Note: 左下方文字是“群作用的不动点集”，中间的虚线箭头即“翻模”。$X$ 与 $X^+$ 是翻模前后的光滑三维簇，按其商掉 $\mu_n$ 后，$X_n$ 带终端奇点而 $X^+_n$ 光滑，故商操作把翻模变成翻转。]

% [原文 p.53]
计算标准除子与例外曲线的交点数并不困难。计算结果如下：
\[
C_n\cdot K_{X_n}=-\frac{n-1}{n},\qquad C_n^+\cdot K_{X_n^+}=n-1.
\]
因此，当 $n\ge 2$ 时，$X^+_n\to Y_n$ 是 $X_n\to Y_n$ 的翻转。

关于此例所具有的性质之一，请注意以下一点。在 $X$ 的孤立固定点处选取坐标系 $(x,v',u)$，则曲线 $C$ 成为 $v'$-轴。$\sigma=(v'^{n-1}-x)(\rd x\wedge\rd v'\wedge\rd u)^{-1}$ 给出 $K_X^{-1}$ 的一个 $\mu_n$-不变局部截面，其零集与曲线 $C$ 的局部交点数等于 $(n-1)$。此截面是 $\mu_n$-不变的，故给出 $K_{X_n}^{-1}$ 的局部截面 $\sigma_n$。记零集为 $D_n=(\sigma_n=0)$。由构造，$D_n\cong\{(v',u)\text{-平面}\}/\mu_n$，它是 $A_{n-1}$ 型的 Du Val 奇点（参见 §4.3）。由于 $C_n\cdot D_n=C_n\cdot K_{X_n}^{-1}$，容易看出 $D_n$ 在 $C_n$ 的Euclid 邻域上属于 $|K_{X_n}^{-1}|$。
\end{example}

上面例子中出现的操作可归纳为如下形式。

\begin{definition}{2.8}
设 $f\colon X\to Y$ 为固有双有理态射，且例外集 $\Exc(f)$ 在 $X$ 内的余维数至少为 2。再设 $K_X$ 为 $\Q$-Cartier ，$-K_X$ 为 $f$-丰富。若簇 $X^+$ 与固有双有理态射 $f^+\colon X^+\to Y$ 组成的对满足下列条件，则称之为 $f$ 的翻转。
\begin{enumerate}[label=（\arabic*）,leftmargin=2.6em,itemsep=1pt]
\item $K_{X^+}$ 为 $\Q$-Cartier ；
\item $K_{X^+}$ 为 $f^+$-丰富；
\item 例外集 $\Exc(f^+)$ 在 $X^+$ 内的余维数至少为 2。
\end{enumerate}

稍有几处记号滥用，双有理映射 $\varphi\colon X\dashrightarrow X^+$ 也称为翻转。已知翻转若存在则唯一，因此问题在于其存在性。翻转给出如下交换图式：

\begin{center}
\begin{tikzcd}[column sep=8em, row sep=3em]
X \arrow[rr, dashed, "\varphi"] \arrow[dr, "f"'] & & X^{+} \arrow[dl, "f^{+}"] \\
& Y &
\end{tikzcd}
\end{center}
其中 $-K_X$ 为 $f$-丰富，$K_{X^{+}}$ 为 $f^{+}$-丰富。
\end{definition}
% [Tutor Note: 翻转与翻模容易混淆。翻转要求 $K_{X^+}$ 为 $f^+$-丰富，即标准类在翻转前后的符号相反；翻模则保持 $K$ 的符号不变，只是把收缩掉的小簇换成另一个小簇。]

% [原文 p.54]
进行翻转之后，$X^+$ 是否比 $X$“更简单”并不显然。不过，上例中 $X$ 带有奇点而 $X^+$ 非奇异，确实变简单了。

\begin{remark}{2.9}
文献中对术语的用法并不统一。高维双有理几何的多数文献遵循定义 2.8，但在许多论文中，只要 $\Exc(f)$ 与 $\Exc(f^+)$ 的余维数都不小于 2，如上那样的任意图式也称为翻转，而对标准类不作任何假设。更一般的翻转概念将在定义 3.33 中引入。\hfill$\square$
\end{remark}

\par\medskip\noindent\textbf{2.10}\quad 在整理上述讨论的形式之前，必须先确定森纲领实际奏效的奇点类。这样的类早在 19 世纪就在如下设定下被研究过。

设 $Y_d\subset\PP^n$ 为带奇点的超曲面，$f\colon X\to Y$ 为其奇点消解。由伴随公式，预期成立下式。
\[
H^0(X,\cO_X(K_X))\overset{?}{=}H^0(Y,\cO_Y(K_Y))=H^0(Y,\cO_Y(d-n-1))=\binom{d-1}{n}.
\]
然而，由于 $Y$ 的奇点强加“伴随条件”，在 $\overset{?}{=}$ 处有时不成立等号，而成立“$<$”。因此，理解哪些奇点不施加伴随条件（即等号成立）是件有趣的事。按近代记号，这些奇点称为标准奇点。\hfill$\square$
% [Tutor Note: 原文此处作“同伴公式”，即通行的“伴随公式”（adjunction formula），下文的“同伴条件”即“伴随条件”。此处只译其义，不另立术语。]

\begin{definition}{2.11}
设 $Y$ 为正规簇。若下列条件满足，则称 $Y$ 具有标准奇点（canonical singularity），或称 $Y$ 的奇点是标准的（canonical）。
\begin{enumerate}[label=（\arabic*）,leftmargin=2.6em,itemsep=1pt]
\item $K_Y$ 为 $\Q$-Cartier ；
\item 对任意奇点消解 $f\colon X\to Y$，$f_*\cO_X(mK_X)=\cO_Y(mK_Y)$ 成立。
\end{enumerate}
（容易看出，条件 (2) 只需对一个奇点消解验证即可。）
\end{definition}

注意下面一点，可定义更小的类。若 $Y$ 非奇异、$f\colon X\to Y$ 为双有理态射、$\gamma\in H^0(Y,\cO_Y(mK_Y))$，则 $f^*\gamma$ 沿例外除子 $\Exc(f)$ 具有零点。

% [原文 p.55]
这比只要求 $f^*\gamma$ 沿 $\Exc(f)$ 没有极点的定义 2.11(2) 更强。这一条件给出如下定义。可以证明，这是纲领能够奏效的最小类。

\begin{definition}{2.12}
设 $Y$ 为正规簇。若下列条件满足，则称 $Y$ 具有终端奇点（terminal singularity），或称 $Y$ 的奇点是终端的（terminal）。
\begin{enumerate}[label=（\arabic*）,leftmargin=2.6em,itemsep=1pt]
\item $K_Y$ 为 $\Q$-Cartier ；
\item 对任意奇点消解 $f\colon X\to Y$，$f_*\cO_X(mK_X-E)=\cO_Y(mK_Y)$ 成立；其中 $E\subset X$ 为把所有例外除子以系数 1 相加所得的除子。
\end{enumerate}
（容易看出，条件 (2) 只需对一个奇点消解验证即可。）
\end{definition}
% [Tutor Note: 标准与终端的差别只在条件 (2)：$f_*\cO_X(mK_X)$ 与 $f_*\cO_X(mK_X-E)$ 之间差一个 $E$。终端要求 $f^*\gamma$ 沿例外除子有零点，即容许的极点阶数更小，故终端类比标准类更小、奇点更轻。]

仿照定义 1.30，引入高维极小模型的概念。

\begin{definition}{2.13}
设 $X$ 为正规且固有的簇。若下列两个条件成立，则称 $X$ 是极小的（minimal），或称 $X$ 为极小模型（minimal model）。
\begin{enumerate}[label=（\arabic*）,leftmargin=2.6em,itemsep=1pt]
\item $X$ 只有终端奇点；
\item $K_X$ 为 nef。
\end{enumerate}
若 $Y$ 为与 $X$ 双有理等价的非奇异固有簇，则也称 $X$ 为 $Y$ 的极小模型。不过，极小模型 $X$ 并非由 $Y$ 唯一决定（参见后文的定义 3.50、定理 3.52(2)）。当 $Y$ 为非奇异射影曲面时，定义 1.30 意义下 $Y$ 的极小模型，按本定义显然也是 $Y$ 的极小模型。逆命题也可经由定义 3.50 证明，但以后用定义 1.30，故省略。
\end{definition}

有了这些定义，就可以逐步精确地叙述森纲领的概要。

\par\medskip\noindent\textbf{2.14}（森纲领，或极小模型纲领）\quad 从非奇异射影簇 $X$ 出发，反复进行易于理解的双有理变换，直到得到整体结构更简单的簇（记为 $X^*$）。以下详列必要的手续。

% [原文 p.56]
\par\medskip\noindent\textbf{步骤 0（初始数据）}\quad 从仅有 $\Q$-分解的（$\Q$-factorial）终端奇点的 $\C$ 上的射影簇 $X=X_0$ 出发。定义如下递归过程：依次构造中间簇 $X_i$，最后以簇 $X^*$ 结束。设已构造到 $X_i$。

\par\medskip\noindent\textbf{步骤 1（准备）}\quad 若 $K_{X_i}$ 为 nef，则无需做什么，直接进入步骤 3 (2)。若 $K_{X_i}$ 不为 nef，则先建立两个结果。
\begin{enumerate}[label=（\arabic*）,leftmargin=2.6em,itemsep=1pt]
\item （锥定理）$\NE(X_i)=\NE(X_i)_{K_{X_i}\ge 0}+\sum\R^+[C_i]$。若 $K_{X_i}$ 不为 nef，则选取 $K_{X_i}$-负的极射线 $R_i\subset\NE(X_i)$。（适当选取 $R_i$ 有时非常有用，但这里试图建立无论怎样选取都能奏效的理论。）
\item （极射线的收缩）用 $\mathrm{cont}_{R_i}\colon X_i\to Y_i$ 表示极射线 $R_i$ 的收缩态射（命题 2.5）。
\end{enumerate}

\par\medskip\noindent\textbf{步骤 2（双有理变换）}\quad 由命题 2.5，可能有 3 种类型的收缩 $\mathrm{cont}_{R_i}\colon X_i\to Y_i$。其中 2 种按如下方式用来构造新的簇 $X_{i+1}$。
\begin{enumerate}[label=（\arabic*）,leftmargin=2.6em,itemsep=1pt]
\item （除子收缩）设 $\mathrm{cont}_{R_i}\colon X_i\to Y_i$ 为命题 2.5(2) 中的除子收缩。此时定义 $X_{i+1}=Y_i$。后面将证明 $X_{i+1}$ 仍只有 $\Q$-分解的终端奇点。承认这一点，就可以对簇 $X_{i+1}$ 从步骤 0 重新开始。
\item （翻转收缩）设 $\mathrm{cont}_{R_i}\colon X_i\to Y_i$ 为命题 2.5(3) 中的翻转收缩。此时定义 $X_{i+1}=X_i^+$，即 $\mathrm{cont}_{R_i}\colon X_i\to Y_i$ 的翻转。后面将证明 $X_{i+1}$ 仍只有 $\Q$-分解的终端奇点。因此，可以对簇 $X_{i+1}$ 从步骤 0 重新开始。
\end{enumerate}

\par\medskip\noindent\textbf{步骤 3（最终结果）}\quad 期待这一过程终将停止，并得到下面两种可能性之一。
\begin{enumerate}[label=（\arabic*）,leftmargin=2.6em,itemsep=1pt]
\item （Fano 纤维空间）设 $\mathrm{cont}_{R_i}\colon X_i\to Y_i$ 为命题 2.5(1) 中的 Fano 收缩。此时定义 $X^*:=X_i$。底空间 $Y_i$ 与纤维的维数都比 $X$ 低，因此通过研究它们，可望获得研究 $X^*$ 的线索。
\end{enumerate}
% [注记：上句起点在 p.56 末段，句尾落在 p.57 首行；按 BRIEF §4 规则 1，整句由本块写完，故此处已含 p.57 首行内容，ch02b 不重复。]
% [Tutor Note: 步骤 2 (2) 的翻转若恒存在，纲领才能继续。三维情形由森于 1988 年证明翻转存在；四维以上至今未完全解决，这正是纲领的瓶颈所在。]
